
\subsubsection{6.1 Key Findings}

The direction of DELTA\_ECE is architecture-specific. Code-specialized pre-training (DeepSeek-Coder) produces positive DELTA\_ECE. Code-adapted fine-tuning (CodeLlama) inverts it. General-purpose pre-training (Llama3) produces insensitivity. The assumption that P(True) from hard problems is less reliable holds only for code-specialized architectures.

Global temperature scaling is insufficient to correct architecture-dependent calibration patterns. CodeLlama's extreme T* = 3.95 reflects pathological global confidence inflation; applying this scaling worsens DELTA\_ECE rather than correcting it.

\subsubsection{6.2 Mechanistic Interpretation}

The architecture-dependent DELTA\_ECE pattern is consistent with a training data composition hypothesis:

\begin{itemize}
\item DeepSeek-Coder (code-specialized pre-training): Hard problems trigger genuine uncertainty in P(True); easy problems trigger appropriate confidence. This produces the expected positive DELTA\_ECE.
\end{itemize}

\begin{itemize}
\item Llama3-8B (general-purpose): Code problems are outside the primary competence domain. Both tiers appear similarly uncertain, producing uniform miscalibration insensitive to difficulty.
\end{itemize}

\begin{itemize}
\item CodeLlama-7B (code fine-tuning on Llama): Fine-tuning on large code repositories may create pattern-matching overconfidence on common Python utilities similar to MBPP easy problems. The model assigns high P(True) to solutions resembling frequently-seen code regardless of correctness, producing ECE(easy) > ECE(hard).
\end{itemize}

This interpretation is exploratory (N=1 per architecture category). Direct confirmation would require analyzing whether CodeLlama's high P(True) on easy-tier problems correlates with similarity to its training corpus.

\subsubsection{6.3 Implications for Practitioners}

\begin{itemize}
\item For code-specialized models (DeepSeek family): confidence signals are less reliable on hard problems; lower thresholds may be appropriate for difficult code.
\item For code-adapted models (CodeLlama family): confidence signals are less reliable on easy-looking problems; overconfidence on common patterns should be expected.
\item For general-purpose models (Llama3 family): difficulty-stratified thresholding provides no benefit; miscalibration is uniform.
\item General recommendation: characterize the model's DELTA\_ECE profile before deploying confidence-based code verification.
\end{itemize}

\subsubsection{6.4 Limitations}

\textbf{L1: k=5 pilot methodology.} With only 6 discrete pass@1 values, tier assignments are coarse. Results are directionally stable across M-sensitivity analysis, but future work should use k >= 20 for finer stratification.

\textbf{L2: CodeLlama easy tier underrepresentation.} The CodeLlama easy tier contains only 37 problems (MBPP-only). The effect size (-0.249) is large relative to expected sampling variance and the CI is entirely negative, but the magnitude should be confirmed with larger-sample replication.

\textbf{L3: Three-model exploratory scope (N=1 per category).} Architecture-specific interpretations are exploratory. Replication with N >= 2 models per category is needed to confirm categorical patterns.

\textbf{L4: Weak P(True)-correctness correlation (r = 0.14-0.20).} P(True) captures a weak but statistically significant confidence signal. DELTA\_ECE estimates partially reflect confidence noise in addition to calibration properties. DeepSeek shows a near-zero correlation (r = -0.046, p = 0.048), indicating its confidence signal is minimally informative at the instance level despite producing architecture-level DELTA\_ECE structure.

\textbf{L5: Base model P(True) validity.} Kadavath et al. (2022) validated P(True) primarily on instruction-tuned models. Base models are not explicitly trained to respond to self-evaluation prompts, and their logprob responses may reflect surface pattern matching rather than genuine self-evaluation. The observed non-degenerate std(c) and significant correlations confirm a non-trivial signal exists, but the interpretation as calibration is less certain than in instruction-tuned settings.
